\documentclass{article}
\usepackage{graphicx} % Required for inserting images
\usepackage{amsmath,amsthm}
\usepackage[english]{babel}
\newtheorem{theorem}{Theorem}[section]
\newtheorem{corollary}{Corollary}[theorem]
\newtheorem{lemma}[theorem]{Lemma}
\title{Lisa de ejercicios - Teoría de grafos}
\author{
	Marcos Ayala \\ 
	\and Juaquin Remon
}
\date{April 2025}

\begin{document}
	
	\maketitle
	
	\section{Ejercicio 39}
	
	
	
	Sea $G$ un árbol con un número finito de vértices. Definimos una función $F(X)$ que toma como entrada un árbol 
	$X$ y devuelve un subgrafo $X'$, donde:
	$$X' = X-leaves(X)$$
	mientras que $|V(X)| > 2$, caso contrario, devolverá X sin modificaciones. 
	
	Si aplicamos $F$ de forma iterativa sobre $G$ y sus sucesivos subgrafos, obtendremos finalmente un subgrafo $G'$ de $G$ tal que $|V(G')| \leq 2$ . Este subgrafo $G'$ representa el núcleo del árbol original y contiene a los vértices que minimizan la excentricidad, es decir, los centros del árbol.
	
	\section{Ejercicio 47}
	\begin{lemma}\label{circuito}
		Sea $C = \{v_0, e_1, v_1, e_2, ..., e_k, v_k = v_0\}$ un circuito, entonces
		$\forall v \in V(C), \quad d_C(v) \equiv 0 \mod 2. $
	\end{lemma}
	Demostrando \ref{circuito}. Sea un vértice arbitrario $v \in V(C)$ y sus apariciones en $C$
	$$v_{i_1} = v, v_{i_2} = v, ... v_{i_r} = v,$$  
	notemos que por cada aparición de $v$, $v_x$ en el circuito, le corresponde dos aristas $e_x$ y $e_{x+1}$ a exepción cuando $x = k$, el cual su $e_{x+1} = e_1$. Por lo tanto $d_C(v) = 2 \cdot r$.
	
	Primero, supongamos que $G$ es euleriano, entonces existe un circuito $C$ que recorre todas las aristas en $E(G)$. Por lo tanto, podemos construir una partición trivial de las aristas $P = \{C\}$, el cual es una partición de las aristas en circuitos.
	
	Segundo, supongamos que $E(G)$ puede ser particionado en circuitos $P = \{C_1, C_2,..., C_k\}$ y $G$ es conexo. Notemos que para cualquier $v \in V(G)$, $d_G(v) = \sum_{C_i \in P} d_{C_i}(v)$. Entonces, usando \ref{circuito} sabemos que $d_{C_i} \equiv 0 \mod 2$, asi que 
	$$d_G(v) \equiv 0 \mod 2$$
	Además, 
	$$d_G(v) \geq 2$$
	ya que si $d_G(v) = 0$ entonces no ha aparecido en ningún circuito, por lo que $P$ no seria una partición de $G$, contradicción.
	Por lo tanto, como para cualquier $v \in V(G)$, $d_G(v) >= 2$ y par, entonces $G$ es un grafo con circuito euleriano.
	
	\section{Ejercicio 59}
	Supongamos por contradicción que $V(P) \cap V(Q) \cap V(R) = \emptyset$. 
	Sea $$V(P) \cap V(Q) = \{a_1, a_2, ...\}$$ $$V(Q) \cap V(R) = \{b_1, b_2, ...\}$$$$V(P) \cap V(R) = C$$
	Podemos construir un camino $S_1$ de $a_1$ a $b_1$ como un subcamino de $Q$, ya que $a_1, b_1 \in V(Q)$.
	
	Sea $s_1 = (a_1, ..., c_i)$ el subcamino de $P$ que minimiza la longitud del camino de $a_i$ a algún vértice de $C$. Note que $c_i$ es el único vértice de $C$ que pertenece a $s_1$, caso contrario no seria el subcamino mínimo. Sea $s_2 = (c_i, ..., b_i)$ el subcamino de $R$. Note que $c_i \notin S_1$, caso contrario, entonces $V(P) \cap V(Q) \cap V(R) \neq \emptyset$. Entonces podemos construir otro camino $S_2 = s_1 + s_2$, $S_2$ es un camino ya que $s_1 \cap s_2 = \{c_i\}$. Entonces pudimos construir dos caminos de $a_1$ a $b_1$, el primero $S_1$ que no contiene ningún vértice de $C$, y el segundo, $S_2$ que contiene a $c_i$. Por lo tanto, no existe un único camino en $T$ de $a_i$ a $b_i$, entonces $T$ no es un árbol, contradicción.
	
	
	\section{Ejercicio 63}
	
	Supongamos por contradicción que una arista puente $e \in E(G)$, donde G es un grafo conexo, no existe en alguno de los árboles generadores $T$ de $G$. Entonces, $T$ no sería conexo, lo cual contradice la definición de árbol generador.
	
	
	
	\section{Ejercicio 73}
	\subsection{G tiene un circuito hamiltoniano}
	Según las condiciones del problema tenemos que los grados de cada conjunto son los siguientes: 
	$$ d(v)_{v \in V_1} = |V_2| + |V_3| = 2|V_1| + 3|V_1|$$
	$$ d(v)_{v \in V_2} = |V_1| + |V_3| = |V_1| + 3|V_1| $$
	$$ d(v)_{v \in V_3} = |V_1| + |V_2| = |V_1| + 2|V_1| $$
	
	Asimismo, también sabemos que el grafo $G$ tiene una cantidad $6|V1|$ de nodos. Por ello, a partir del teorema de Dirac, el cual menciona que cualquier grafo con $n \geq 3 $ y $\delta(G) \geq \lceil \frac{n}{2} \rceil $ tiene un circuito hamiltoniano, podemos decir que G tiene un circuito hamiltoniano, pues $d(v)_{v \in V_1} = 5|V_1| \geq 3|V_1|$, $d(v)_{v \in V_2} = 4|V_1| \geq 3|V_1|$ y $d(v)_{v \in V_3} = 3|V_1| \geq 3|V_1|$ . 
	
	\subsection{G no tiene un circuito hamiltoniano}
	Según las condiciones del problema tenemos que los grados de cada conjunto son los siguientes: 
	$$ d(v)_{v \in V_1} = |V_2| + |V_3| = 2|V_1| + 3|V_1| + 1$$
	$$ d(v)_{v \in V_2} = |V_1| + |V_3| = |V_1| + 3|V_1| +1 $$
	$$ d(v)_{v \in V_3} = |V_1| + |V_2| = |V_1| + 2|V_1| $$
	
	Supongamos, por contradicción, que bajo la condición dada en el problema, el grafo \( G \) contiene un circuito hamiltoniano. Recordemos que, en un grafo bipartito completo, para que exista un ciclo hamiltoniano es necesario que ambos conjuntos de vértices tengan la misma cantidad de elementos.
	
	En nuestro caso, \( G \) es un grafo tripartito completo, cuyos conjuntos de vértices tienen tamaños \( |V_1| \), \( |V_2| = 2|V_1| \) y \( |V_3| = 3|V_1| + 1 \). Observamos que los vértices de \( V_1 \) y \( V_2 \) pueden conectarse completamente con los vértices de \( V_3 \). Sin embargo, debido a que \( V_3 \) tiene un vértice adicional, quedaría un vértice sin emparejar, imposibilitando así la construcción de un circuito hamiltoniano. Esto contradice nuestra suposición inicial, por lo tanto, \( G \) no contiene un circuito hamiltoniano.
	
\end{document}



\end{document}
